\scanpage{8}
\[
\pi_1(\underline{\mathrm{End}})\overset{?}{=}
\mathrm{Aut}\,(u=\mathrm{id},v=\mathrm{id},\alpha\deleted{\unsure{\in}}=1\in G^I)=\ ?
\]
$\lambda\in\mathbb Z^I$, $\mu\in G$ tels que

a) $\mathrm{id}_G=\mathrm{int}(\mu)\circ\mathrm{id}_G$ i.e. $\mu\in Z(G)$ (centre)

b) $\forall i,\quad\mu=h_i^{\lambda_i}$

Or, sauf le cas de la sphère à $1$ ou $2$ trous, le groupe $G$ est libre à $\unsure{N|n}\geq2$ générateurs, donc son centre est nul, de plus les $h_i$ sont libres. Donc dans ce cas on a $\mu=1$, $\lambda_i=0\ \forall i\in I$, donc $\boxed{\pi_1(\underline{\mathrm{End}})=0}$

Dans le cas de la sphère à $2$ trous, \inserted{i.e. le cylindre,} on a $G\simeq\mathbb Z$, $h_i=\pm1$, $\deleted{\unsure{h_j=\pm1}}$ ; les $h_i$ sont encore libres (donc les $\lambda_i$ déterminés par $\mu$) et $G=Z(G)$. Donc
\[
\pi_1(\underline{\mathrm{End}})=G\quad(\simeq\mathbb Z)
\]
\margin{(pas canoniquement)}
Dans le cas de la sphère à $1$ trou, $G=0$, et on trouve $\pi_1(\underline{\mathrm{End}})\simeq\mathbb Z$ (isom. canonique)

\begin{center}\rule{1.5cm}{.4pt}\end{center}

Regardons maintenant la $2$-catégorie des $(C,D,\varphi)$ \uline{avec $D$ fixé}, donc $u=\mathrm{id}$, $\lambda=\mathrm{id}$. On trouve pour $\underline{\mathrm{End}}_{D\text{ fixé}}(C,D,\varphi)$

objets : couples $(v,\alpha)$ avec
\[
\left\{\begin{aligned}
&v\in\mathrm{Aut}\,G,\\
&\alpha\in G^I\\
&v(h_i)=\mathrm{int}(\alpha(i))\cdot h_i
\end{aligned}\right.
\]
$\mathrm{Hom}\big((v_1,\alpha_1),(v_2,\alpha_2)\big)={}$ ens. des $\mu\in G$ tels que
\[
\left\{\begin{aligned}
&v_2=\mathrm{int}(\mu)\circ v_1\\
&\alpha_2(i)=\mu\alpha_1(i)
\end{aligned}\right.
\]
On trouve que $(v_1,\alpha_1)\simeq(v_2,\alpha_2)$ $\deleted{\unsure{\sim}}\Rightarrow\deleted{\unsure{\exists}}$ \deleted{\unsure{$\mu$ tel q}} \inserted{$\widetilde v_1=\widetilde v_2$} i.e $\exists\mu_0$ tel que $v_2=\mathrm{int}(\mu_0)v_1$
\[
\begin{aligned}
\text{d'où}\qquad v_2(\unsure{h^i|h_i|h'_i|\ell_i})&=\underline{\mathrm{int}(\mu_0)}\,\underline{v_1}(h_i)\\
\Vert\qquad &\hspace{5em}\Vert\\
\mathrm{int}(\alpha_2(i))h_i &\qquad\mathrm{int}(\mu_0)\,\mathrm{int}\,\alpha_1(i)\,v_1(h_i)
\end{aligned}
\]
donc $\alpha_2(i)^{-1}\mu_0\alpha_1(i)$ centralise $h_i$
% archive sheet number: 7
